% Meta: ~300 palabras
\section{Uso de la ciencia de datos en la inclusión financiera del sector agro}

Las Reglas de Operación de los apoyos al sector agropecuario buscan ampliar el acceso al financiamiento formal y mejorar la productividad de las unidades rurales \citep{fira2025reglas}. Un obstáculo es la información: quien evalúa un crédito suele conocer la zona mejor que la parcela, y el rendimiento anterior llega cuando la decisión ya se tomó. Un estimado por polígono, construido con datos públicos de satélite y clima, no reemplaza el expediente ni la visita, pero sí adelanta una magnitud en toneladas por hectárea y un intervalo de error conocido.

En la práctica, el dashboard permite tres usos alineados con esa política. El primero es de tamiz: ordenar parcelas de una zona por el estimado y por el residual histórico, para concentrar la asistencia técnica donde el modelo y el clima de la temporada se separan. El segundo es de conversación con el productor: la ficha muestra el número, las variables que más empujan esa parcela y la aclaración de que la explicación no cubre toda la mezcla. El tercero es de límite institucional: el RMSE de 0.508 t/ha donde se califica, y el deterioro en Singuilucan o al salir del municipio, indican hasta dónde el dato puede pesar en una resolución de crédito. Usarlo como si fuera un aforo medido reproduciría la exclusión por otro camino, el de un modelo opaco.

La viabilidad es modesta. No hace falta una red de sensores en la parcela: hacen falta el polígono, el ciclo y no extrapolar a otra región ni a otro año sin reentrenar. El código y los artefactos permiten repetir la consulta. Si el estimado se usara en operación, el siguiente paso sería contrastarlo con las etiquetas ocultas y, sobre todo, con más de una temporada.
